\begin{textsample}{NEG Dim 8 – human – Score -2.00 – mg/mg\_ef\_matematica\_4\_p5.txt}  \label{ex:f8_neg_004}
Na Matemática escolar, o processo de aprender uma noção em um contexto, abstrair e depois aplicá -la em outro contexto envolve capacidades essenciais, como formular, empregar, interpretar e \textbf{avaliar} – criar, enfim –, e não somente a resolução de enunciados típicos que são, muitas vezes, meros exercícios e apenas simulam alguma aprendizagem.
Assim, algumas das habilidades formuladas começam por : “ resolver e \textbf{elaborar} problemas envolvendo... ”.
Nessa enunciação está implícito que se pretende não apenas a resolução do problema, mas também que os estudantes reflitam e questionem o que ocorreria se algum dado do problema fosse alterado ou se alguma condição fosse acrescida ou retirada.
Nessa perspectiva, pretende -se que os estudantes também formulem problemas em outros contextos. ( BNCC, 2017, p.277 ).

% matched lemmas: avaliar, elaborar
\end{textsample}
